%%
%% part2-book-components/ch15-index.tex
%%
%% Chapter 15: Index — xindy with Persian sorting, main
%% entries, sub-entries, and cross-references.
%%

\chapter{نمایه}
\label{chap:index}

\section{چرا این موضوع مهم است؟}
\label{sec:index-why}

نمایه یکی از ابزارهای مهم ناوبری در کتاب‌های
فنی و علمی است. یک نمایه‌ی حرفه‌ای به خواننده
کمک می‌کند که مفاهیم مورد نظر را سریع پیدا
کند و مسیر خود را در کتاب گم نکند.

\lr{MatinBook} از \lr{xindy} با ماژول
\lr{persian-variant2} برای ساخت نمایه استفاده
می‌کند. این ترکیب، مرتب‌سازی صحیح فارسی را
تضمین می‌کند.

\mnote{بسته‌ی \lr{makeindex} نمی‌تواند فارسی را درست مرتب کند و حروف پ، چ، ژ، گ، ک را اشتباه می‌چیند. به همین دلیل \lr{MatinBook} از \lr{xindy} استفاده می‌کند.}

\section{ساختار نمایه}
\label{sec:index-structure}

نمایه از چند نوع مدخل تشکیل شده است:

\begin{itemize}
    \item \textbf{مدخل اصلی:} یک واژه یا عبارت
    که در نمایه ظاهر می‌شود.
    \item \textbf{زیرمدخل:} یک مدخل فرعی که
    زیر مدخل اصلی قرار می‌گیرد.
    \item \textbf{زیرزیرمدخل:} عمیق‌ترین سطح
    مدخل.
    \item \textbf{ارجاع «نگاه کنید به»:} ارجاع
    از یک واژه به واژه‌ی دیگر.
    \item \textbf{ارجاع «همچنین نگاه کنید به»:}
    ارجاع تکمیلی.
\end{itemize}

\mnote{مرتب‌سازی نمایه در \lr{MatinBook} بر اساس الفبای فارسی انجام می‌شود و با \lr{xindy} به‌درستی کار می‌کند.}

\section{مدخل اصلی}
\label{sec:index-main}

برای ایجاد یک مدخل اصلی، از دستور \cmd{index}
استفاده کنید:

\begin{latin}
\begin{minted}{latex}
|\pc{الگوریتم}|\index{|\pc{الگوریتم}|}|\pc{یک مفهوم اساسی است.}|%
\end{minted}
\end{latin}

خروجی: الگوریتم\index{الگوریتم} یک مفهوم اساسی است.

\mnote{دستور \lr{\textbackslash index} واژه را در نمایه ثبت می‌کند ولی در متن چیزی نمایش نمی‌دهد.}

\section{زیرمدخل}
\label{sec:index-sub}

برای ایجاد یک زیرمدخل، از علامت \lr{!} استفاده
کنید:

\begin{latin}
\begin{minted}{latex}
|\pc{برنامه‌نویسی}|\index{|\pc{برنامه‌نویسی!پایتون}|}|\pc{یک زبان محبوب است.}|%
\end{minted}
\end{latin}

خروجی: برنامه‌نویسی\index{برنامه‌نویسی!پایتون} یک
زبان محبوب است.

\mnote{در نمایه، این مدخل به‌صورت «برنامه‌نویسی ← پایتون» نمایش داده می‌شود.}

\subsection{زیرزیرمدخل}
\label{sec:index-subsub}

برای زیرزیرمدخل، دو علامت \lr{!} استفاده
کنید:

\begin{latin}
\begin{minted}{latex}
|\pc{ریاضیات}|\index{|\pc{ریاضیات!جبر!گروه}|}|\pc{شاخه‌ای از ریاضیات است.}|%
\end{minted}
\end{latin}

\mnote{توصیه می‌شود بیش از دو سطح زیرمدخل نداشته باشید، چون نمایه پیچیده می‌شود.}

\section{دستورات کمکی}
\label{sec:index-helpers}

\lr{MatinBook} چند دستور کمکی برای نمایه ارائه
می‌دهد. جدول \cref{tab:index-helpers} این
دستورات را خلاصه می‌کند.

\begin{matintable}{دستورات کمکی نمایه}{tab:index-helpers}
\begin{tblr}{
    width = \textwidth,
    colspec = {l X[l]},
    hlines, vlines,
    colsep = 8pt,
    rowsep = 4pt,
    row{1} = {font=\bfseries},
}
دستور & کاربرد \\
\lr{\textbackslash idxbold\{...\}} & مدخل با فونت سیاه \\
\lr{\textbackslash idxitalic\{...\}} & مدخل با فونت مورب \\
\lr{\textbackslash idxsee\{...\}\{...\}} & ارجاع «نگاه کنید به» \\
\lr{\textbackslash idxseealso\{...\}\{...\}} & ارجاع «همچنین نگاه کنید به» \\
\lr{\textbackslash idxsub\{...\}\{...\}} & زیرمدخل سریع \\
\lr{\textbackslash idxsubsub\{...\}\{...\}\{...\}} & زیرزیرمدخل سریع \\
\end{tblr}
\end{matintable}

\subsection{مدخل با فونت سیاه}
\label{sec:index-bold}

برای نمایش یک مدخل با فونت سیاه، از
\cmd{idxbold} استفاده کنید:

\begin{latin}
\begin{minted}{latex}
|\pc{مفهوم مهم}|\idxbold{|\pc{مفهوم مهم}|}%
\end{minted}
\end{latin}

\mnote{مدخل‌های سیاه برای مفاهیم کلیدی استفاده می‌شوند و در نمایه برجسته‌تر دیده می‌شوند.}

\subsection{ارجاع «نگاه کنید به»}
\label{sec:index-see}

برای ارجاع از یک واژه به واژه‌ی دیگر، از
\cmd{idxsee} استفاده کنید:

\begin{latin}
\begin{minted}{latex}
\idxsee{|\pc{پایتون}|}{|\pc{زبان برنامه‌نویسی}|}%
\end{minted}
\end{latin}

خروجی: در نمایه، این مدخل به‌صورت «پایتون ← نگاه
کنید به زبان برنامه‌نویسی» نمایش داده می‌شود.

\mnote{ارجاع «نگاه کنید به» برای واژه‌های مترادف یا مفاهیم مرتبط استفاده می‌شود.}

\subsection{ارجاع «همچنین نگاه کنید به»}
\label{sec:index-seealso}

برای ارجاع تکمیلی، از \cmd{idxseealso} استفاده
کنید:

\begin{latin}
\begin{minted}{latex}
\idxseealso{|\pc{برنامه‌نویسی}|}{|\pc{الگوریتم}|}%
\end{minted}
\end{latin}

\mnote{ارجاع «همچنین نگاه کنید به» برای مفاهیم مرتبط ولی نه مترادف استفاده می‌شود.}

\subsection{زیرمدخل سریع}
\label{sec:index-subfast}

برای ایجاد زیرمدخل سریع، از \cmd{idxsub}
استفاده کنید:

\begin{latin}
\begin{minted}{latex}
\idxsub{|\pc{ریاضیات}|}{|\pc{آنالیز}|}%
\end{minted}
\end{latin}

خروجی: این دستور معادل
\cmd{index\{ریاضیات!آنالیز\}} است.

\mnote{دستور \lr{\textbackslash idxsub} معادل \lr{\textbackslash index\{اصلی!زیر\}} است و کوتاه‌تر نوشته می‌شود.}

\section{مرتب‌سازی فارسی}
\label{sec:index-sorting}

مرتب‌سازی نمایه در \lr{MatinBook} با \lr{xindy}
و ماژول \lr{persian-variant2} انجام می‌شود.
این ماژول، ترتیب صحیح حروف فارسی را رعایت
می‌کند.

جدول \cref{tab:index-sorting} ترتیب حروف
مشکل‌دار در فارسی را نشان می‌دهد.

\begin{matintable}{ترتیب حروف مشکل‌دار در فارسی}{tab:index-sorting}
\begin{tblr}{
    width = \textwidth,
    colspec = {c X[l] X[l]},
    hlines, vlines,
    colsep = 8pt,
    rowsep = 4pt,
    row{1} = {font=\bfseries},
}
حرف & ترتیب صحیح & ترتیب \lr{makeindex} \\
پ & بعد از ب & اشتباه \\
چ & بعد از ج & اشتباه \\
ژ & بعد از ز & اشتباه \\
گ & بعد از ک & اشتباه \\
ک & بعد از گ & اشتباه \\
\end{tblr}
\end{matintable}

\mnote{اگر از \lr{makeindex} استفاده کنید، این حروف در جای اشتباه قرار می‌گیرند و نمایه نامرتب می‌شود.}

\section{چاپ نمایه}
\label{sec:index-print}

برای چاپ نمایه، از دستور \cmd{printindex} در
\lr{backmatter} استفاده کنید:

\begin{latin}
\begin{minted}{latex}
\backmatter
\frontmattergeometry

\printbibliography[title={|\pc{منابع و مراجع}|}]

\printindex
\end{minted}
\end{latin}

\mnote{گزینه‌ی \lr{intoc} در \lr{mb-index.sty} باعث می‌شود نمایه در فهرست مطالب ظاهر شود.}

\section{چرخه‌ی کامپایل}
\label{sec:index-compile}

برای ساخت نمایه، باید چرخه‌ی کامل کامپایل را
اجرا کنید:

\begin{latin}
\begin{minted}{bash}
# |\pc{مرحله‌ی اول: کامپایل اول}|%
xelatex -shell-escape main.tex

# |\pc{مرحله‌ی دوم: پردازش نمایه}|%
# |\pc{\lr{xindy} به‌طور خودکار توسط \lr{imakeidx} اجرا می‌شود}|%

# |\pc{مرحله‌ی سوم: کامپایل دوم}|%
xelatex -shell-escape main.tex

# |\pc{مرحله‌ی چهارم: کامپایل سوم}|%
xelatex -shell-escape main.tex
\end{minted}
\end{latin}

\mnote{اگر از \lr{-shell-escape} استفاده نکنید، \lr{xindy} اجرا نمی‌شود و نمایه خالی می‌ماند.}

\section{تنظیمات نمایه}
\label{sec:index-settings}

تنظیمات نمایه در فایل \file{mb-index.sty} تعریف
شده است:

\begin{latin}
\begin{minted}{latex}
\makeindex[
    intoc,
    program=xindy,
    options={-L persian-variant2 -C utf8
             -M texindy -M page-ranges}
]
\end{minted}
\end{latin}

\mnote{گزینه‌ی \lr{intoc} نمایه را به فهرست مطالب اضافه می‌کند. \lr{persian-variant2} نیز مرتب‌سازی صحیح فارسی را تضمین می‌کند.}

\section{مثال کامل}
\label{sec:index-example}

در ادامه، یک مثال کامل از استفاده‌ی نمایه در
یک فصل آورده شده است:

\begin{latin}
\begin{minted}{latex}
\chapter{|\pc{مفاهیم برنامه‌نویسی}|}

\section{|\pc{زبان‌های برنامه‌نویسی}|}

|\pc{پایتون}|\index{|\pc{پایتون}|}|\pc{یک زبان سطح بالا است.}|%

|\pc{برنامه‌نویسی شیءگرا}|\index{|\pc{برنامه‌نویسی!شیءگرا}|}%
|\pc{یک پارادایم است.}|%

\idxsee{|\pc{جاوااسکریپت}|}{|\pc{زبان برنامه‌نویسی وب}|}%

\idxseealso{|\pc{الگوریتم}|}{|\pc{ساختمان داده}|}%
\end{minted}
\end{latin}

\mnote{هر چه مدخل‌های بیشتری در کتاب داشته باشید، نمایه مفصل‌تر و کاربردی‌تر خواهد بود.}

\section{خطاهای رایج}
\label{sec:index-mistakes}

\subsection{استفاده از \lr{makeindex} به‌جای \lr{xindy}}
\label{sec:index-mistake-makeindex}

\textbf{مشکل:} اگر از \lr{makeindex} استفاده
کنید، حروف پ، چ، ژ، گ، ک اشتباه مرتب می‌شوند.

\textbf{راه‌حل:} همیشه از \lr{xindy} استفاده
کنید. این تنظیم در \file{mb-index.sty} از پیش
انجام شده است.

\subsection{فراموش کردن \lr{-shell-escape}}
\label{sec:index-mistake-shellescape}

\textbf{مشکل:} اگر \lr{-shell-escape} را
فراموش کنید، \lr{xindy} اجرا نمی‌شود و نمایه
خالی می‌ماند.

\textbf{راه‌حل:} همیشه از
\lr{xelatex -shell-escape} استفاده کنید.

\subsection{مدخل‌های تکراری}
\label{sec:index-mistake-duplicate}

\textbf{مشکل:} اگر یک واژه را چند بار با
\cmd{index} علامت‌گذاری کنید، در نمایه چند
بار ظاهر می‌شود.

\textbf{راه‌حل:} واژه را فقط در محل اصلی
تعریف کنید و در جاهای دیگر به آن ارجاع دهید.

\subsection{مدخل‌های بسیار زیاد}
\label{sec:index-mistake-toomany}

\textbf{مشکل:} اگر مدخل‌های نمایه بسیار زیاد
باشند، نمایه از حد معمول بزرگ‌تر می‌شود و
کاربردی‌بودن خود را از دست می‌دهد.

\textbf{راه‌حل:} فقط مفاهیم کلیدی را در نمایه
بگنجانید. قاعده‌ی تجربی: ۵ تا ۱۰ مدخل در هر
صفحه.

\section{تمرین‌ها}
\label{sec:index-exercises}

\begin{exercise}
    یک فصل ساده با سه مدخل اصلی و دو زیرمدخل
    بنویسید و نمایه را چاپ کنید.
    \label{exc:index-basic}
\end{exercise}

\begin{exercise}
    یک مدخل با فونت سیاه و یک مدخل با فونت
    مورب بسازید.
    \label{exc:index-bold}
\end{exercise}

\begin{exercise}
    از \cmd{idxsee} و \cmd{idxseealso} برای
    ارجاع بین مدخل‌ها استفاده کنید.
    \label{exc:index-see}
\end{exercise}

\begin{exercise}
    نمایه‌ی خود را با نمایه‌ی یک کتاب فارسی
    دیگر مقایسه کنید. آیا ترتیب حروف
    یکسان است؟
    \label{exc:index-compare}
\end{exercise}

\section{خلاصه}
\label{sec:index-summary}

در این فصل، با نمایه در \lr{MatinBook} آشنا
شدیم:

\begin{itemize}
    \item \lr{MatinBook} از \lr{xindy} با
    ماژول \lr{persian-variant2} برای مرتب‌سازی
    صحیح فارسی استفاده می‌کند.

    \item \lr{makeindex} نمی‌تواند فارسی را
    درست مرتب کند و باید از \lr{xindy} استفاده
    کرد.

    \item مدخل اصلی با \cmd{index\{...\}}.

    \item زیرمدخل با \cmd{index\{اصلی!زیر\}}.

    \item زیرزیرمدخل با
    \cmd{index\{اصلی!زیر!زیر\}}.

    \item دستورات کمکی: \cmd{idxbold}،
    \cmd{idxitalic}، \cmd{idxsee}،
    \cmd{idxseealso}، \cmd{idxsub}،
    \cmd{idxsubsub}.

    \item نمایه با \cmd{printindex} در
    \lr{backmatter} چاپ می‌شود.

    \item خطاهای رایج شامل استفاده از
    \lr{makeindex}، فراموش کردن
    \lr{-shell-escape}، مدخل‌های تکراری، و
    مدخل‌های بسیار زیاد است.
\end{itemize}

در فصل بعدی (\cref{chap:cross-refs})، با
ارجاع هوشمند در \lr{MatinBook} آشنا می‌شوید
و یاد می‌گیرید که چگونه به عناصر مختلف کتاب
ارجاع دهید.